Il \autoref{cap:capitoli/04c-forward-backward} ha descritto come si calcolano il forward e la backpropagation di una proiezione STERA. Questo capitolo descrive il software che li esegue: il motore di training, scritto in Python con PyTorch per GPU NVIDIA, e i motori di inferenza, scritti in Python con MLX e kernel Metal per la GPU del chip Apple M1. Per ciascuno descriviamo l'algoritmo di decompressione delle coordinate, i kernel principali e il costo delle operazioni. Il codice è nel repository del progetto: il motore di training in \texttt{training/}, i motori di inferenza in \texttt{systems/engines/}.

\section{Il motore di training}
\label{sec:motore-training}

\subsection{Linguaggio e piattaforma}
Il motore di training è scritto in Python con PyTorch ed è eseguito con CUDA su GPU NVIDIA; STERA 30M è stato addestrato su due GPU P100. I processi, uno per GPU, sono avviati con \texttt{torchrun}: ciascuno elabora una parte del batch, i gradienti dei due processi sono sommati con un'operazione di riduzione collettiva (NCCL) e ogni processo applica lo stesso aggiornamento. Per limitare la memoria, il batch di ogni GPU è elaborato in gruppi di otto sequenze, accumulando i gradienti prima dell'aggiornamento. Tutti i calcoli sono in FP32; la precisione ridotta TF32 delle GPU è disattivata.

\subsection{Decompressione delle coordinate}
\label{sec:ricostruzione-training}
A ogni passo il motore ricostruisce le matrici con un algoritmo a tabelle. Una volta sola, all'avvio, calcola tre tabelle di \num{16384} elementi, le stesse per tutti i blocchi: per ogni casella $(i,j)$ del blocco, l'indice della coordinata di diagonale o di base che vi compare, l'indice dell'eventuale correzione e un indicatore che vale 1 se la correzione c'è e 0 altrimenti. La ricostruzione di tutti i blocchi di una matrice si riduce allora a due letture indicizzate del vettore delle coordinate, una somma e una moltiplicazione per la scala:
\[
\Phi(c)_{ij}=c_{\mathrm{base}(i,j)}+\mathrm{maschera}(i,j)\,c_{\mathrm{corr}(i,j)} ,
\]
seguite dalla moltiplicazione per i guadagni e dalla somma di $UV^\top$. Sono tutte operazioni standard di PyTorch, eseguite in parallelo su tutti i blocchi.

Il vantaggio di questa scelta riguarda la backpropagation. L'autograd di PyTorch conosce la derivata di ogni operazione usata: il passo all'indietro di una lettura indicizzata è una somma indicizzata, che accumula il gradiente di ogni casella nella coordinata letta. Il gradiente delle coordinate si ottiene quindi senza codice dedicato, ed è esattamente quello della \autoref{sec:gradiente-coordinate}. Il resto del passo usa i prodotti densi della GPU, l'operazione più ottimizzata delle librerie numeriche.

\subsection{Inizializzazione delle coordinate}
\label{sec:init-coordinate}
Per le matrici dense si estrae di solito ogni peso da una normale di media nulla e deviazione standard $\sigma_0$ piccola, per esempio $0{,}02$. In un blocco RowMask23 gli elementi fuori diagonale con una correzione sono la somma di \emph{due} coordinate indipendenti: se ogni coordinata ha deviazione standard $\sigma$, la varianza media di un elemento vale $\sigma^2\big(1 + \num{2921}/\num{16384}\big)$. Per avere la stessa varianza media di una matrice densa si sceglie quindi
\begin{equation}
\sigma = \frac{\sigma_0}{\sqrt{1+\num{2921}/\num{16384}}}.
\end{equation}


\subsection{Disposizione delle coordinate in memoria}
\label{sec:layout-coordinate}
Le coordinate di un blocco sono memorizzate in un unico vettore di lunghezza \num{3176}: nelle posizioni $0,\dots,127$ la diagonale ($d_i$ in posizione $i$), nelle posizioni $128,\dots,254$ la base ($b_s$ in posizione $127+s$), nelle posizioni $255,\dots,3175$ le correzioni ($e_{i,k}$ in posizione $255+23(i-1)+k$). Ogni coordinata ha così un indice calcolabile in tempo costante, sia nella ricostruzione sia nell'accumulo del gradiente.


\section{I motori di inferenza}
\label{sec:motori-inferenza}

\subsection{Linguaggio e piattaforma}
I motori di inferenza sono scritti in Python con MLX, la libreria di calcolo di Apple per i chip della serie M. Le operazioni che non esistono in MLX sono scritte come kernel nel linguaggio Metal della GPU e compilate a tempo di esecuzione attraverso l'interfaccia \texttt{mx.fast.metal\_kernel} di MLX. I motori sono due: il motore denso, che espande una volta tutte le matrici e usa i prodotti di MLX, e il motore compatto, che applica le proiezioni direttamente alle coordinate. Entrambi leggono gli stessi pesi FP32 del checkpoint finale.

All'avvio il motore compatto riordina le correzioni di ogni blocco: la correzione $e_{i,k}$ passa dalla posizione $255+23(i-1)+k$ del training alla posizione $255+127k+(i-1)$. Il riordino non cambia nessun valore; serve perché i 127 thread che elaborano le righe di un blocco leggano, per ogni $k$, correzioni contigue in memoria.

\subsection{Kernel di espansione}
\label{sec:kernel-rowmask}
Il primo kernel espande le coordinate di una matrice nella matrice densa corrispondente. Ogni thread della GPU calcola un elemento: dalla sua posizione $(r,c)$ ricava il blocco a cui appartiene, la posizione $(i,j)$ dentro il blocco e l'indice $s=i\oplus j$; legge quindi la diagonale, la base ed eventualmente la correzione dal vettore delle coordinate del blocco e scrive il risultato moltiplicato per la scala e, nella versione usata nel prefill, per i guadagni. Per sapere se una posizione ha una correzione, e quale, il kernel usa una tabella di $128\times128$ elementi, la stessa per tutti i blocchi. Il risultato è identico, bit per bit, a quello della ricostruzione del motore di training; il controllo è eseguito automaticamente su ogni matrice prima delle misure.

\subsection{Kernel del prodotto diretto}
Il secondo kernel calcola il prodotto $y=Wx$ senza costruire $W$, con i quattro passi della \autoref{sec:forward-stera} e la formula \eqref{eq:wht-conv} per la parte XOR:
\begin{enumerate}
\item l'ingresso è moltiplicato per $g_{\mathrm{in}}$ e di ogni suo segmento di 128 componenti si calcola una volta la trasformata di Walsh--Hadamard, condivisa da tutti i blocchi della stessa colonna;
\item la GPU lancia un gruppo di 128 thread per ogni blocco di righe, e ciascun gruppo accumula in registro i contributi di più blocchi della stessa riga, così da non scrivere in memoria un risultato parziale per ogni blocco. Per ogni blocco, il thread della riga $i$ calcola il contributo della diagonale e delle 23 correzioni, leggendo l'ingresso agli indici $i\oplus\sigma(i,k)$, e il contributo della base, moltiplicando la trasformata dell'ingresso per $Hb$ e applicando la trasformata inversa. La trasformata inversa usa i primi cinque stadi all'interno di gruppi di 32 thread, con scambi diretti fra registri, e gli ultimi due attraverso la memoria condivisa del gruppo;
\item il risultato è moltiplicato per $g_{\mathrm{out}}$ e vi si somma il termine low-rank, calcolato come $U(V^\top x)$ con i prodotti di MLX.
\end{enumerate}
Il vettore $Hb$ di ogni blocco dipende solo dalle coordinate: il motore lo calcola all'avvio e lo tiene in memoria, 128 valori per blocco, cioè \num{633600} valori e $2{,}53$~MB per STERA 30M, il $2{,}1\%$ dei pesi. Questa tabella non è salvata su disco e non fa parte dei byte dei pesi riportati nel \autoref{cap:capitoli/07-sistemi}, ma è compresa nella memoria del processo misurata. La correttezza del prodotto diretto è verificata confrontandolo con il prodotto ottenuto espandendo la matrice; sui prompt usati nelle misure, i motori compatto e denso producono la stessa sequenza di token.

\subsection{Prefill con espansione al bisogno e low-rank a 8 bit}
Nel prefill il motore compatto può usare l'espansione al bisogno della \autoref{sec:guida-inferenza-stera}: per ogni layer espande con il kernel di espansione la parte strutturata delle sue sei matrici, esegue i prodotti densi fra matrici per tutti i token del prompt e rilascia le matrici prima di passare al layer successivo, così che in memoria resti espanso un solo layer alla volta; il termine low-rank non viene espanso. Nella variante con il low-rank a 8 bit, i fattori $U$ e $V$ sono quantizzati con una scala e uno spostamento ogni 32 valori e il prodotto $U(V^\top x)$ usa il prodotto quantizzato di MLX, senza riconvertire i fattori in FP32.

\section{Costo delle operazioni}
\label{sec:costi-motore}

La \autoref{tab:costi-motore}, subito sotto, riassume il costo delle operazioni principali per una proiezione $m\times n$ con $PQ=mn/16384$ blocchi, rango $r$ e $T$ token, contando le operazioni aritmetiche e i valori dei pesi da leggere.

\begin{table}[H]
\centering
\small
\renewcommand{\arraystretch}{1.25}
\begin{tabular}{>{\raggedright\arraybackslash}p{40mm}>{\raggedright\arraybackslash}p{50mm}>{\raggedright\arraybackslash}p{38mm}}
\toprule
Operazione & Operazioni aritmetiche & Valori dei pesi letti \\
\midrule
Prodotto denso & $T\,mn$ moltiplicazioni e somme & $mn$ \\
\rowcolor{steraBlu!10} Prodotto diretto sulle coordinate & $T\,\big(PQ\,(\num{3049}+N+N\log_2N)+Q\,N\log_2N\big)$, cioè circa $0{,}25\,T\,mn$ & $\num{3176}\,PQ+N\,PQ\approx0{,}2\,mn$ \\
Espansione di una matrice & $mn$ scritture & $\num{3176}\,PQ\approx0{,}19\,mn$ \\
Termine low-rank $U(V^\top x)$ & $T\,r(m+n)$ & $r(m+n)$ \\
\rowcolor{steraArancio!12} Backpropagation nel training & ricostruzione ($mn$), due prodotti densi ($2\,T\,mn$), somma indicizzata ($mn$) & $mn$ ricostruiti \\
\bottomrule
\end{tabular}
\caption{Costo delle operazioni principali su una proiezione STERA $m\times n$, con $N=128$, $N\log_2N=896$ e $T$ token. Il prodotto diretto conta, per ogni blocco, diagonale e correzioni ($128+\num{2921}=\num{3049}$ moltiplicazioni e somme), il prodotto con $Hb$ ($N$) e la trasformata inversa ($N\log_2N$), più la trasformata dell'ingresso una volta per colonna di blocchi; i valori letti sono le coordinate e il vettore $Hb$ di ogni blocco.}
\label{tab:costi-motore}
\end{table}

Nel decode, con $T=1$, il prodotto diretto esegue circa un quarto delle operazioni del prodotto denso e legge circa un quinto dei valori: poiché la generazione token per token è limitata dalla lettura dei pesi (\autoref{sec:memory-bound}), è questa seconda riduzione a determinare il vantaggio. Nel prefill, con $T$ grande, il costo è dominato dal calcolo, e il prodotto denso fra matrici, che riusa ogni peso per tutti i token, diventa più conveniente: per questo il prefill usa l'espansione al bisogno, il cui costo $mn$ si ripartisce su tutti i token del prompt. Nel training il motore sceglie il percorso denso per tutte le operazioni, perché con batch di migliaia di token il costo della ricostruzione è piccolo rispetto ai prodotti.

\section{Sintesi}

Il training usa PyTorch su GPU NVIDIA: ricostruisce le matrici con tabelle di indici e lascia all'autograd la somma dei gradienti nelle coordinate. L'inferenza usa MLX e due kernel Metal: l'espansione delle coordinate, identica bit per bit alla ricostruzione del training, e il prodotto diretto, che legge circa un quinto dei valori del prodotto denso. Le scelte fra percorso denso e prodotto diretto seguono il costo delle operazioni: denso quando molti token condividono gli stessi pesi, diretto quando se ne elabora uno alla volta.
